\section{Some Plane Geometry}

Let $\gamma\colon[a,b]\to\mathbb R^2$ be admissible. It is a \emph{simple
closed curve} if $\gamma(a)=\gamma(b)$ and $\gamma$ is injective on $[a,b)$.
On each smooth segment, its unit tangent is
\[
T(t)=\frac{\gamma'(t)}{|\gamma'(t)|}\in S^1.
\]

If $\gamma$ is continuously differentiable, a \emph{tangent angle function}
is a continuous $\theta\colon[a,b]\to\mathbb R$ satisfying
\[
T(t)=(\cos\theta(t),\sin\theta(t)).
\]
It exists by lifting $T$ through the covering
$q(s)=(\cos s,\sin s)$.
\uses{prop:lifting-properties-covering}
The unique lifting property implies that any two tangent angle functions differ
by a constant in $2\pi\mathbb Z$.
\uses{prop:lifting-properties-covering}
If in addition $\gamma$ is simple closed and
$\gamma'(a)=\gamma'(b)$, its \emph{rotation index} is
\[
\rho(\gamma)=\frac{\theta(b)-\theta(a)}{2\pi}\in\mathbb Z.
\]
This integer is independent of the tangent angle function.

Now let $(a_0,\ldots,a_k)$ be an admissible partition for an admissible simple
closed curve, with $a_0=a$ and $a_k=b$. The points $\gamma(a_i)$ are its
\emph{vertices}, and the restrictions to $[a_{i-1},a_i]$ are its \emph{edges}.
At a vertex, write $T(a_i^-)$ and $T(a_i^+)$ for the left and right unit
tangents. A vertex is
\begin{itemize}
\item \emph{ordinary} if $T(a_i^-)\ne\pm T(a_i^+)$;
\item \emph{flat} if $T(a_i^-)=T(a_i^+)$;
\item a \emph{cusp} if $T(a_i^-)=-T(a_i^+)$.
\end{itemize}

At an ordinary vertex, the \emph{exterior angle} $\varepsilon_i$ is the
oriented angle from $T(a_i^-)$ to $T(a_i^+)$ in $(-\pi,\pi)$. It is positive
exactly when the ordered pair $(T(a_i^-),T(a_i^+))$ is positively oriented.
At a flat vertex set $\varepsilon_i=0$; at a cusp it is undefined because the
choices $\pi$ and $-\pi$ cannot be distinguished. At the common endpoint
$\gamma(a)=\gamma(b)$, use $T(b)$ as the left tangent and $T(a)$ as the right
tangent. For an ordinary or flat vertex, the \emph{interior angle} is
\[
\vartheta_i=\pi-\varepsilon_i,
\qquad 0<\vartheta_i<2\pi.
\]

A \emph{curved polygon} is an admissible simple closed plane curve without
cusp vertices whose image is the boundary of a precompact open set
$\Omega\subseteq\mathbb R^2$. It is \emph{positively oriented} when its smooth
tangent agrees with the boundary orientation induced by Stokes's theorem;
equivalently, $\Omega$ lies to its left.

A tangent angle function for a curved polygon is a piecewise continuous
$\theta\colon[a,b]\to\mathbb R$ that represents $T$, is right-continuous at
each $a_i$, and incorporates the exterior-angle jumps
\[
\theta(a_i)=\lim_{t\nearrow a_i}\theta(t)+\varepsilon_i
\qquad (9.1)
\]
at the interior vertices, together with
\[
\theta(b)=\lim_{t\nearrow b}\theta(t)+\varepsilon_k.
\qquad (9.2)
\]
Successive path lifting on the edges gives such a function, unique up to a
constant in $2\pi\mathbb Z$. Define
\[
\rho(\gamma)=\frac{\theta(b)-\theta(a)}{2\pi}\in\mathbb Z.
\]

\begin{theorem}[Rotation Index Theorem]
\label{thm:rotation-index-theorem}
\dcref{ch9:9.1}
Every positively oriented curved polygon in the plane has rotation index
$+1$.
\end{theorem}

\begin{proof}
\sketch First suppose all vertices are flat. Reparametrize by arclength; the
one-sided unit tangents then agree at every vertex, so the periodically
extended curve is continuously differentiable. Write its period as $b-a$. If
$\theta\colon\mathbb R\to\mathbb R$ lifts its periodic unit tangent, uniqueness
of lifts gives
\begin{equation}
\theta(t+b-a)=\theta(t)+2\pi\rho(\gamma). \tag{9.3}
\end{equation}
Thus every restriction to an interval of length $b-a$ has the same rotation
index. Choose the initial parameter at a minimum of the $y$-coordinate,
translate that point to the origin, and use positive orientation to obtain
\[
T(a)=T(b)=(1,0),
\qquad \theta(a)=0.
\]
The curve lies in the closed upper half-plane.

On
\[
\Delta=\{(t_1,t_2):a\leq t_1\leq t_2\leq b\},
\]
define the secant-direction map
\[
V(t_1,t_2)=
\begin{cases}
\dfrac{\gamma(t_2)-\gamma(t_1)}
{|\gamma(t_2)-\gamma(t_1)|},
&t_1<t_2\text{ and }(t_1,t_2)\ne(a,b),\\[2mm]
T(t_1),&t_1=t_2,\\
-T(b),&(t_1,t_2)=(a,b).
\end{cases}
\]
The fundamental theorem of calculus gives, uniformly near the diagonal,
\[
\frac{\gamma(t_2)-\gamma(t_1)}{t_2-t_1}\longrightarrow\gamma'(t)
\quad\text{as }(t_1,t_2)\to(t,t),
\]
so $V$ is continuous there. Periodicity gives continuity at $(a,b)$. Since
$\Delta$ is simply connected, $V$ has a unique lift
$\varphi\colon\Delta\to\mathbb R$ with $\varphi(a,a)=0$. Along the side
$t_1=a$, secants point into the closed upper half-plane, so
$\varphi(a,t_2)\in[0,\pi]$ and $\varphi(a,b)=\pi$. Along the side $t_2=b$,
they point into the closed lower half-plane, so
$\varphi(t_1,b)\in[\pi,2\pi]$ and $\varphi(b,b)=2\pi$. On the diagonal,
$\varphi$ is a tangent angle function; hence
\[
\rho(\gamma)
=\frac{\varphi(b,b)-\varphi(a,a)}{2\pi}=1.
\]

For ordinary vertices, choose the initial endpoint away from them and round
each corner without changing the rotation index. At a vertex $\gamma(a_i)$
with exterior angle $\varepsilon_i$, choose
\[
0<\alpha<\frac12(\pi-|\varepsilon_i|).
\]
The jump convention gives
\[
\lim_{t\nearrow a_i}\theta(t)=\theta(a_i)-\varepsilon_i.
\]
Choose $\delta>0$ so that
\[
|\theta(t)-(\theta(a_i)-\varepsilon_i)|<\alpha
\quad(a_i-\delta<t<a_i),
\]
and
\[
|\theta(t)-\theta(a_i)|<\alpha
\quad(a_i<t<a_i+\delta).
\]
Isolate this portion of the curve in a ball about $\gamma(a_i)$, and let
$t_1<a_i<t_2$ be its entry and exit times. The total tangent-angle change
$\Delta\theta$ on $[t_1,t_2]$ satisfies
\[
\varepsilon_i-2\alpha<\Delta\theta<\varepsilon_i+2\alpha,
\qquad -\pi<\Delta\theta<\pi.
\]
Replace this piece by a smooth segment with the same endpoints and endpoint
velocities whose tangent angle varies monotonically from $\theta(t_1)$ to
$\theta(t_2)$. Its angle change is the unique representative in
$(-\pi,\pi)$, hence equals $\Delta\theta$. Repeating at every ordinary vertex
produces a positively oriented simple closed curve with continuous velocity
and the same rotation index, so the first case applies.
\end{proof}
\section{The Gauss--Bonnet Formula}

Let $(M,g)$ be an oriented Riemannian $2$-manifold.  An admissible simple
closed curve $\gamma:[a,b]\to M$ is a \emph{curved polygon} if its image is
the boundary of a precompact open set $\Omega\subseteq M$ and
$\overline\Omega$ lies in an oriented coordinate disk that sends $\gamma$ to
a planar curved polygon.  The set $\Omega$ is the \emph{interior}; if every
edge is a geodesic segment, $\gamma$ is a \emph{geodesic polygon}.

The curve is \emph{positively oriented} when its parametrization has the
Stokes boundary orientation.  On each smooth edge, set
\[
T(t)=\frac{\gamma'(t)}{|\gamma'(t)|_g}.
\]
At an ordinary vertex $\gamma(a_i)$, the \emph{exterior angle} is the oriented
angle from $T(a_i^-)$ to $T(a_i^+)$,
\begin{equation}
\varepsilon_i=
\frac{dV_g(T(a_i^-),T(a_i^+))}
{|dV_g(T(a_i^-),T(a_i^+))|}
\arccos\langle T(a_i^-),T(a_i^+)\rangle_g. \tag{9.4}
\end{equation}
At a flat vertex set $\varepsilon_i=0$.  The interior angle is
$\theta_i=\pi-\varepsilon_i$; the endpoint vertex $\gamma(a)=\gamma(b)$ is
treated identically.

Choose an oriented chart containing $\overline\Omega$, and apply
Gram--Schmidt to its coordinate frame to obtain an oriented orthonormal frame
$(E_1,E_2)$, with $E_1$ a positive multiple of the first coordinate vector.
A \emph{tangent angle function} is a piecewise continuous
$\theta:[a,b]\to\mathbb R$, right-continuous at vertices and satisfying the
prescribed exterior-angle jumps, such that
\[
T(t)=\cos\theta(t)E_1|_{\gamma(t)}
+\sin\theta(t)E_2|_{\gamma(t)}.
\]
It exists because
\begin{equation}
T(t)=u_1(t)E_1|_{\gamma(t)}+u_2(t)E_2|_{\gamma(t)}, \tag{9.5}
\end{equation}
where the piecewise continuous map $(u_1,u_2)$ takes values in $S^1$ and
lifts locally to an angle.  Define the \emph{rotation index} by
\[
\rho(\gamma)=\frac{\theta(b)-\theta(a)}{2\pi}.
\]

\begin{lemma}[Rotation Index of a Positive Curved Polygon]
\label{lem:rotation-index-plus-one}
\uses{thm:rotation-index-theorem}
\dcref{ch9:9.2}

Every positively oriented curved polygon in an oriented Riemannian
$2$-manifold has rotation index $+1$.
\end{lemma}

\begin{proof}
\sketch
Transfer the curve and metric to the coordinate plane, and let $\bar g$ be the
Euclidean metric there.  For $s\in[0,1]$, set
$g_s=sg+(1-s)\bar g$.  The rotation index $\rho_{g_s}(\gamma)$ is integral for
every $s$.  The Gram--Schmidt frame varies continuously with $s$; on each
smooth edge the coefficients
\[
u_j(t,s)=\langle T_s(t),E_j^{(s)}|_{\gamma(t)}\rangle_{g_s},
\qquad
T_s=\frac{\gamma'}{|\gamma'|_{g_s}},
\]
have a continuous angle lift, and (9.4) makes every vertex angle continuous
in $s$.  Thus $s\mapsto\rho_{g_s}(\gamma)$ is continuous and integer-valued,
so it is constant.  The planar rotation-index theorem gives
$\rho_g(\gamma)=\rho_{\bar g}(\gamma)=1$.
\end{proof}

Henceforth parametrize $\gamma$ by unit speed.  On each smooth edge, let $N$
be the unique unit normal for which $(\gamma',N)$ is positively oriented.  If
$\gamma$ is positively oriented, $N$ points inward along $\partial\Omega$.
The \emph{signed geodesic curvature} is
\[
\kappa_N=\langle D_T\gamma',N\rangle_g.
\]
Since $|\gamma'|=1$,
$D_T\gamma'=\kappa_NN$ and the unsigned geodesic curvature is
$|\kappa_N|$.  The sign is positive when the curve bends toward $\Omega$.

\begin{theorem}[The Gauss--Bonnet Formula]
\label{thm:gauss-bonnet-formula}
\uses{prop:arc-length-integral, thm:rotation-index-theorem, lem:rotation-index-plus-one, prob:arc-length-integral}
\dcref{ch9:9.3}

Let $\gamma$ be a positively oriented curved polygon in an oriented
Riemannian $2$-manifold, with interior $\Omega$ and exterior angles
$\varepsilon_1,\ldots,\varepsilon_k$.  Then
\begin{equation}
\int_\Omega K\,dA+\int_\gamma\kappa_N\,ds
+\sum_{i=1}^k\varepsilon_i=2\pi, \tag{9.6}
\end{equation}
where $K$ is Gaussian curvature, $dA$ is the oriented Riemannian area form,
and the boundary integral is with respect to arc length.
\end{theorem}

\begin{proof}
\sketch
Let $(a_0,\ldots,a_k)$ be an admissible partition, choose oriented
coordinates on a neighborhood of $\overline\Omega$, and let $\theta$ be a
tangent angle function.  The rotation-index lemma and the fundamental theorem
of calculus give
\begin{equation}
2\pi=\theta(b)-\theta(a)
=\sum_{i=1}^k\varepsilon_i
+\sum_{i=1}^k\int_{a_{i-1}}^{a_i}\theta'(t)\,dt. \tag{9.7}
\end{equation}

For the oriented orthonormal frame $(E_1,E_2)$ used above,
\[
\gamma'=\cos\theta E_1+\sin\theta E_2,
\qquad
N=-\sin\theta E_1+\cos\theta E_2.
\]
Differentiation along each smooth edge gives
\begin{equation}
D_t\gamma'=\theta'N+(\cos\theta)\nabla_{\gamma'}E_1
+(\sin\theta)\nabla_{\gamma'}E_2. \tag{9.8}
\end{equation}
Define the connection $1$-form
\[
\omega(v)=\langle E_1,\nabla_vE_2\rangle
=-\langle\nabla_vE_1,E_2\rangle.
\]
Metric compatibility yields
\begin{align}
\nabla_vE_1&=-\omega(v)E_2,\\
\nabla_vE_2&=\omega(v)E_1.
\label{eq:covariant-frame-derivatives}
\tag{9.9}
\end{align}
Substituting these identities into (9.8) gives
\[
\kappa_N=\langle D_t\gamma',N\rangle
=\theta'-\omega(\gamma').
\]
Equation (9.7) becomes
\[
2\pi=\sum_{i=1}^k\varepsilon_i
+\int_\gamma\kappa_N\,ds+\int_\gamma\omega.
\]
It remains to prove
\begin{equation}
\int_\gamma\omega=\int_\Omega K\,dA. \tag{9.10}
\end{equation}

The closure $\overline\Omega$ is a smooth manifold with corners, so Stokes's
theorem identifies the left side with $\int_\Omega d\omega$.  Since
$dA(E_1,E_2)=1$, the curvature convention and (9.9) give
\begin{align*}
K
&=\mathrm{Rm}(E_1,E_2,E_2,E_1)\\
&=\left\langle
\nabla_{E_1}\nabla_{E_2}E_2
-\nabla_{E_2}\nabla_{E_1}E_2
-\nabla_{[E_1,E_2]}E_2,E_1\right\rangle\\
&=E_1(\omega(E_2))-E_2(\omega(E_1))-\omega([E_1,E_2])\\
&=d\omega(E_1,E_2).
\end{align*}
Thus $d\omega=K\,dA$, proving (9.10) and (9.6).
\end{proof}

\begin{corollary}[Angle-Sum Theorem]
\label{cor:angle-sum-theorem}
\uses{thm:gauss-bonnet-formula}
\dcref{ch9:9.4}

The interior angles of a Euclidean triangle sum to $\pi$.
\end{corollary}

\begin{corollary}[Circumference Theorem]
\label{cor:circumference-theorem}
\uses{thm:gauss-bonnet-formula}
\dcref{ch9:9.5}

A Euclidean circle of radius $R$ has circumference $2\pi R$.
\end{corollary}

\begin{corollary}[Total Curvature Theorem]
\label{cor:total-curvature-theorem}
\uses{thm:gauss-bonnet-formula}
\dcref{ch9:9.6}

If $\gamma:[a,b]\to\mathbb R^2$ is smooth, unit speed, simple, and closed,
$\gamma'(a)=\gamma'(b)$, and $N$ is its inward normal, then
\[
\int_a^b\kappa_N(t)\,dt=2\pi.
\]
\end{corollary}
\section{The Gauss--Bonnet Theorem}

A \emph{curved triangle} in a smooth compact surface is a curved polygon with
three edges and three vertices. A \emph{smooth triangulation} is a finite
collection of curved triangles whose interiors are pairwise disjoint, whose
closures cover the surface, and such that two distinct triangles intersect, if
at all, in one common vertex or one common edge. Every smooth compact surface
admits such a triangulation. For a triangulated surface, define
\[
\chi(M)=V-E+F,
\]
where $V,E,F$ are its numbers of vertices, edges, and faces. This number is
independent of the triangulation.

\begin{theorem}[The Gauss--Bonnet Theorem]
\label{thm:gauss-bonnet-theorem}
\uses{thm:gauss-bonnet-formula}
\dcref{ch9:9.7}

If $(M,g)$ is a smoothly triangulated compact Riemannian $2$-manifold without
boundary, then
\[
\int_MK\,dA=2\pi\chi(M),
\]
where $K$ is the Gaussian curvature and $dA$ is the Riemannian density.
\end{theorem}

\begin{proof}
\sketch It suffices to treat each connected component. First suppose $M$ is
connected and oriented, and orient every face boundary by the induced boundary
orientation. Write the faces as $\Omega_i$, their edges as $\gamma_{ij}$, and
their interior angles as $\theta_{ij}$, with
$1\leq i\leq F$ and $1\leq j\leq3$. Summing the Gauss--Bonnet formula over all
faces gives
\begin{equation}
\sum_{i=1}^F\int_{\Omega_i}K\,dA
+\sum_{i=1}^F\sum_{j=1}^3\int_{\gamma_{ij}}\kappa_N\,ds
+\sum_{i=1}^F\sum_{j=1}^3(\pi-\theta_{ij})
=2\pi F. \tag{9.11}
\end{equation}
Every edge occurs twice with opposite boundary orientations, so the geodesic
curvature integrals cancel. Hence
\begin{equation}
\int_MK\,dA+3\pi F-
\sum_{i=1}^F\sum_{j=1}^3\theta_{ij}=2\pi F. \tag{9.12}
\end{equation}
The incident angles at each vertex sum to $2\pi$, so the double angle sum is
$2\pi V$, and therefore
\begin{equation}
\int_MK\,dA=2\pi V-\pi F. \tag{9.13}
\end{equation}
Counting face-edge incidences gives $3F=2E$, and consequently
\[
2\pi V-\pi F
=2\pi V-2\pi E+2\pi F
=2\pi\chi(M).
\]

Now suppose $M$ is nonorientable. Let
$\pi\colon\widetilde M\to M$ be its connected orientable two-sheeted cover.
It is compact, and for $\widetilde g=\pi^*g$,
\[
\widetilde K=\pi^*K,
\qquad d\widetilde A=\pi^*dA,
\qquad
\int_{\widetilde M}\widetilde K\,d\widetilde A
=2\int_MK\,dA.
\]
The triangulation lifts to $\widetilde M$. Indeed, if $\Omega$ is the interior
of a curved triangle, choose a disk chart whose domain contains
$\overline\Omega$ and which identifies it with the closure
$\overline\Omega_0$ of a plane curved triangle. The resulting embedding
\[
F\colon\overline\Omega_0\longrightarrow\overline\Omega
\]
has exactly two disjoint lifts because the disk is simply connected and the
cover has two sheets. Applying this to every triangle gives a triangulation of
$\widetilde M$ with
\[
\widetilde V=2V,
\qquad \widetilde E=2E,
\qquad \widetilde F=2F,
\qquad \chi(\widetilde M)=2\chi(M).
\]
Apply the oriented case to $\widetilde M$ and divide by $2$.
\end{proof}

For a compact connected orientable surface of genus $n$,
$\chi=2-2n$; for a compact connected nonorientable surface of genus $n$,
$\chi=2-n$.

\begin{corollary}
\label{cor:gaussian-curvature-and-topology}
\uses{thm:gauss-bonnet-theorem}
\dcref{ch9:9.8}

Let $(M,g)$ be a compact Riemannian $2$-manifold and let $K$ be its Gaussian
curvature.
\begin{enumerate}
\item[(a)] If $M$ is homeomorphic to the sphere or projective plane, then
$K>0$ somewhere.
\item[(b)] If $M$ is homeomorphic to the torus or Klein bottle, then either
$K\equiv0$ or $K$ takes both positive and negative values.
\item[(c)] On every other compact surface, $K<0$ somewhere.
\end{enumerate}
\end{corollary}

\begin{corollary}
\label{cor:curvature-determines-cover}
\uses{cor:gaussian-curvature-and-topology, prop:compact-covering-finite-sheeted, prop:fiber-cardinality-fundamental-group, exer:A.62}
\dcref{ch9:9.9}

Let $(M,g)$ be a compact Riemannian $2$-manifold with Gaussian curvature $K$.
\begin{enumerate}
\item[(a)] If $K>0$ everywhere, then the universal covering manifold is
homeomorphic to $S^2$, and $\pi_1(M)$ is trivial or isomorphic to
$\mathbb Z/2$.
\item[(b)] If $K\leq0$ everywhere, then the universal covering manifold is
homeomorphic to $\mathbb R^2$, and $\pi_1(M)$ is infinite.
\end{enumerate}
\end{corollary}

\begin{proof}
\sketch If $K>0$, Gauss--Bonnet gives $\chi(M)>0$. Classification of compact
surfaces leaves only $S^2$ and $\mathbb{RP}^2$, whose universal cover is
$S^2$ and whose fundamental groups are respectively trivial and
$\mathbb Z/2$.

If $K\leq0$, then $\chi(M)\leq0$. Thus $M$ is orientable of genus at least
$1$ or nonorientable of genus at least $2$. The universal cover is
$\mathbb R^2$ for the torus and Klein bottle and $\mathbb H^2$ otherwise; in
all cases it is homeomorphic to $\mathbb R^2$ and is noncompact. A finite-sheeted
cover of compact $M$ would be compact, so the universal covering has infinitely
many sheets. Its fiber cardinality equals $|\pi_1(M)|$, which is therefore
infinite.
\end{proof}

For a compact manifold of any dimension, the Euler characteristic is a
triangulation-independent topological invariant; it vanishes in odd dimensions
and may be nonzero in even dimensions (see \cite[Thm. 13.36]{LeeTM}). For each
oriented $2n$-dimensional inner product space $V$, the
\emph{Chern--Gauss--Bonnet theorem} provides a basis-independent Pfaffian map
\[
\operatorname{Pf}\colon\mathcal R(V^*)\longrightarrow\Lambda^{2n}(V^*)
\]
such that every oriented compact Riemannian $2n$-manifold satisfies
\[
\int_M\operatorname{Pf}(\operatorname{Rm})=(2\pi)^n\chi(M).
\]
The normalization of $\operatorname{Pf}$ determines the constant on the right;
see \cite[Vol. 5]{Spi79}. In dimension four, the convention used here gives
\begin{equation}
\int_M\left(
|\operatorname{Rm}|^2-4|\operatorname{Rc}|^2+S^2
\right)\,dV_g
=32\pi^2\chi(M).
\label{eq:chern-gauss-bonnet-4d}
\tag{9.14}
\end{equation}

The Hopf conjecture asserts that a compact even-dimensional manifold admitting
a metric of strictly positive sectional curvature has positive Euler
characteristic. It holds in dimensions $2$ and $4$ and remains open in higher
dimensions; see \cite[p. 320]{Pet16}.

\begin{exercise}[Angle excess formula]
\label{prob:angle-excess-formula}
\dcref{ch9:9-2}
\uses{thm:gauss-bonnet-formula}
Let $(M, g)$ be a Riemannian 2-manifold. If $\gamma$ is a geodesic polygon in $M$ with $n$ vertices, the \textbf{angle excess} of $\gamma$ is defined as
\[
E(\gamma) = \left(\sum_{i=1}^{n} \theta_i\right) - (n-2)\pi,
\]
where $\theta_1, \ldots, \theta_n$ are the interior angles of $\gamma$. Show that if $M$ has constant Gaussian curvature $K$, then every geodesic polygon has angle excess equal to $K$ times the area of the region bounded by the polygon.
\end{exercise}

\begin{exercise}
\label{prob:compact-lie-groups-isomorphic}
\dcref{ch9:9-8}
\uses{prob:lie-group-bi-invariant-metric,thm:gauss-bonnet-theorem}
Use the Gauss--Bonnet theorem to prove that every compact connected Lie
group of dimension 2 is isomorphic to the direct product group $S^1 \times S^1$.
[Hint: See Problem 8-17.]
\end{exercise}

\begin{exercise}
\label{prob:cylinder-metric-gaussian-curvature}
\dcref{ch9:9-11}
Suppose $g$ is a Riemannian metric on the cylinder $S^1 \times [0,1]$ such that both
boundary curves are totally geodesic. Prove that the Gaussian curvature of
$g$ either is identically zero or attains both positive and negative values. Give
examples of both possibilities.
\end{exercise}

\begin{exercise}
\label{prob:einstein-manifold-euler-characteristic}
\dcref{ch9:9-13}
\uses{eq:chern-gauss-bonnet-4d}
Use the four-dimensional Chern--Gauss--Bonnet formula (9.14) to prove that
a compact 4-dimensional Einstein manifold must have positive Euler characteristic unless it is flat.
\end{exercise}

\begin{exercise}[Embedded submanifold cannot be nonpositively curved]
\label{prob:embedded-submanifold-positive-curvature}
\dcref{ch9:9-6}
Let $M \subseteq \mathbb{R}^3$ be a compact, embedded, 2-dimensional Riemannian submanifold. Show that $M$ cannot have $K \leq 0$ everywhere. [Hint: Look at a point where the distance from the origin takes a maximum.]
\end{exercise}

\begin{exercise}[Gauss--Bonnet formula on a sphere]
\label{prob:gauss-bonnet-sphere}
\dcref{ch9:9-3}
\uses{thm:gauss-bonnet-formula}
Given $h \in (-R, R)$, let $C_h$ be the circle in $S^2(R) \subseteq \mathbb{R}^3$ where $z = h$ (where we label the standard coordinates in $\mathbb{R}^3$ as $(x, y, z)$), and let $\Omega$ be the subset of $S^2(R)$ where $z > h$. Compute the signed curvature of $C_h$ and verify the Gauss--Bonnet formula in this case.
\end{exercise}

\begin{exercise}[Gauss--Bonnet theorem on the torus]
\label{prob:gauss-bonnet-torus}
\dcref{ch9:9-4}
\uses{thm:gauss-bonnet-theorem}
Let $T \subseteq \mathbb{R}^3$ be the torus of revolution obtained by revolving the circle $(r - 2)^2 + z^2 = 1$ around the $z$-axis (see p.\ 19). Compute the Gaussian curvature of $T$ and verify the Gauss--Bonnet theorem in this case.
\end{exercise}

\begin{exercise}[The Gauss--Bonnet theorem for surfaces with boundary]
\label{prob:gauss-bonnet-with-boundary}
\dcref{ch9:9-10}
\uses{thm:gauss-bonnet-formula,thm:gauss-bonnet-theorem}
Suppose $(M, g)$ is a compact Riemannian 2-manifold with boundary, endowed
with a smooth triangulation such that the intersection of each curved triangle
with $\partial M$, if not empty, is either a single vertex or a single edge. Then
\[
\int_M K \, dA + \int_{\partial M} \kappa_N \, ds = 2\pi \chi(M),
\]
where $\kappa_N$ is the signed geodesic curvature of $\partial M$ with respect to the inward-pointing normal $N$.
\end{exercise}

\begin{exercise}[Geodesic polygons with nonpositive curvature]
\label{prob:geodesic-polygons-negative-curvature}
\dcref{ch9:9-1}
\uses{thm:gauss-bonnet-formula}
Let $(M, g)$ be an oriented Riemannian 2-manifold with nonpositive Gaussian curvature everywhere. Prove that there are no geodesic polygons with exactly 0, 1, or 2 ordinary vertices. Give examples of all three if the curvature hypothesis is not satisfied.
\end{exercise}

\begin{exercise}
\label{prob:ideal-triangles-hyperbolic}
\dcref{ch9:9-9}
\uses{thm:gauss-bonnet-formula}
(a) Show that there is an upper bound for the areas of geodesic triangles in
the hyperbolic plane $\mathbb{H}^2(R)$, and compute the least upper bound.

(b) Two distinct maximal geodesics in the hyperbolic plane $\mathbb{H}^2$ are said to be
\textit{asymptotically parallel} if they have unit-speed parametrizations $\gamma_1, \gamma_2:
\mathbb{R} \to \mathbb{H}^2$ such that $d_g(\gamma_1(t), \gamma_2(t))$ remains bounded as $t \to +\infty$ or as
$t \to -\infty$. An \textit{ideal triangle} in $\mathbb{H}^2$ is a region whose boundary consists of
three distinct maximal geodesics, any two of which are asymptotically parallel to each other. Show that all ideal triangles have the same finite
area, and compute it. Be careful to justify any limits.
\end{exercise}

\begin{exercise}
\label{prob:plane-curve-classification}
\dcref{ch9:9-12}
\uses{thm:plane-curve-classification}
Prove the plane curve classification theorem (Theorem 1.5). [Hint: Show that
every smooth unit-speed plane curve $\gamma(t) = (x(t), y(t))$ satisfies the second-order ODE $\gamma''(t) = \kappa_N(t) N(t)$, where $N$ is the unit normal vector field given
by $N(t) = (-y'(t), x'(t))$.] (Used on p. 4.)
\end{exercise}

\begin{exercise}
\label{prob:similar-triangles-congruent}
\dcref{ch9:9-7}
Suppose $M$ is either the 2-sphere of radius $R$ or the hyperbolic plane of
radius $R$ for some $R > 0$. Show that similar triangles in $M$ are congruent.
More precisely, if $\gamma_1$ and $\gamma_2$ are geodesic triangles in $M$ such that corresponding side lengths are proportional and corresponding interior angles are
equal, then there exists an isometry of $M$ taking $\gamma_1$ to $\gamma_2$.
\end{exercise}

\begin{exercise}[Smooth triangulation of 2-manifolds]
\label{prob:triangulation-2manifold}
\dcref{ch9:9-5}
This problem outlines a proof that every compact smooth 2-manifold has a smooth triangulation.
\begin{enumerate}
\item[(a)] Show that it suffices to prove that there exist finitely many convex geodesic polygons whose interiors cover $M$, and each of which lies in a uniformly normal convex geodesic ball. (A curved polygon is called convex if the union of the polygon and its interior is a geodesically convex subset of $M$.)
\item[(b)] Using Theorem 6.17, show that there exist finitely many points $v_1, \ldots, v_k$ and a positive number $\varepsilon$ such that the geodesic balls $B_{3\varepsilon}(v_i)$ are geodesically convex and uniformly normal, and the balls $B_\varepsilon(v_i)$ cover $M$.
\item[(c)] For each $i$, show that there is a convex geodesic polygon in $B_{3\varepsilon}(v_i)$ whose interior contains $B_\varepsilon(v_i)$. [Hint: Let the vertices be sufficiently nearby points on the circle of radius $2\varepsilon$ around $v_i$.]
\item[(d)] Prove the result.
\end{enumerate}
(Used on p.\ 276.)
\end{exercise}
